\chapter{Coalgebre}

\begin{definition}[Coalgebra]
  Data una categoria \(\mathcal{C}\) e un endofuntore \(F: \mathcal{C} \to \mathcal{C}\):
  \begin{enumerate}
    \item Una \(F\)-\textbf{coalgebra} (si può omettere il funtore \(F\) se non crea confusione) consiste di un oggetto \(X\in\mathcal{C}\) e un morfismo \(c:X \to F(X)\). \(X\) prende il nome di \textit{spazio degli stati} o \textit{insieme portante} della coalgebra, mentre \(c\) è detta \textit{transizione} o \textit{struttura} della coalgebra.
    \item Un \textbf{omomorfismo di coalgebre} da una coalgebra \(c: X \to F(X)\) a \(d: Y \to F(Y)\) consiste di un morfismo \(f: X \to Y\) che rende il diagramma seguente commutativo:
          \[
            \begin{tikzcd}[column sep=large, row sep=large]
              F(X) \arrow[r, "F(f)"] & F(Y) \\
              X \arrow[u, "c"] \arrow[r, "f"'] & Y \arrow[u, "d"']
            \end{tikzcd}
          \]
    \item Le \(F\)-coalgebre con i loro omomorfismi formano una categoria \(\mathbf{CoAlg}(F)\), che ammette un funtore \(\mathbf{CoAlg}(F) \to \mathcal{C}\) che associa a ogni coalgebra \(c: X \to F(X)\) il suo spazio degli stati \(X\), e a ogni omomorfismo \(f\) lo stesso morfismo visto in \(\mathcal{C}\).
  \end{enumerate}
\end{definition}

La commutatività del diagramma degli omomorfismi esprime una condizione di compatibilità: effettuare una transizione in \(X\) e poi applicare \(f\) produce lo stesso risultato che applicare \(f\) e poi effettuare una transizione in \(Y\). Un omomorfismo di coalgebre è quindi una mappa che rispetta la dinamica di entrambi i sistemi.

L'endofuntore \(F\) è il parametro fondamentale della teoria: esso determina il \textit{tipo di interfaccia} del sistema, cioè la struttura degli stati successori e delle osservazioni accessibili. Per esempio, il funtore \(\mathrm{Seq}_A(X) = \{\bot\} \sqcup (A \times X)\) descrive sistemi che possono terminare o produrre un elemento di \(A\) con un successore \cite{jacobs2017}.

\section{Coalgebra finale e coricorsione}

La nozione di oggetto finale, introdotta nel capitolo precedente, si specializza nella categoria \(\mathbf{CoAlg}(F)\) alla nozione di \textit{coalgebra finale}.

\begin{definition}[Coalgebra finale]
  Una \(F\)-\textbf{coalgebra finale} è un oggetto finale nella categoria \(\mathbf{CoAlg}(F)\), ovvero una coalgebra \(\zeta : Z \to F(Z)\) tale che per ogni coalgebra \(c : X \to F(X)\) esiste un unico omomorfismo di coalgebre \(\mathrm{beh}_c : X \to Z\):
  \begin{center}
    \begin{tikzcd}[column sep=large, row sep=large]
      F(X) \arrow[r, "F(\mathrm{beh}_c)", dashed] & F(Z) \\
      X \arrow[u, "c"] \arrow[r, "\mathrm{beh}_c"', dashed] & Z \arrow[u, "\zeta"']
    \end{tikzcd}
  \end{center}
\end{definition}

Il morfismo \(\mathrm{beh}_c\) è detto \textit{funzione di comportamento}: essa mappa ogni stato \(x \in X\) al suo \textit{comportamento osservabile} come elemento di \(Z\). Poiché la coalgebra finale è un oggetto finale in \(\mathbf{CoAlg}(F)\), dalla proposizione del capitolo precedente segue che essa è unica a meno di isomorfismo unico.

La struttura di transizione di una coalgebra finale gode di una proprietà notevole: è essa stessa un isomorfismo.

\begin{lemma}[Punto fisso di Lambek]
  Se \(\zeta : Z \to F(Z)\) è una coalgebra finale per l'endofuntore \(F\), allora \(\zeta\) è un isomorfismo, e in particolare \(Z \cong F(Z)\).
\end{lemma}

\begin{proof}
  Applicando \(F\) alla coalgebra finale \(\zeta : Z \to F(Z)\) si ottiene una nuova coalgebra \(F(\zeta) : F(Z) \to F(F(Z))\). Per finalità esiste un unico omomorfismo \(f : F(Z) \to Z\) tale che il seguente diagramma commuti:
  \begin{center}
    \begin{tikzcd}[column sep=large, row sep=large]
      F(F(Z)) \arrow[r, "F(f)"] & F(Z) \\
      F(Z) \arrow[u, "F(\zeta)"] \arrow[r, "f"'] & Z \arrow[u, "\zeta"']
    \end{tikzcd}
  \end{center}
  Mostriamo che \(f = \zeta^{-1}\). La composizione \(f \circ \zeta : Z \to Z\) è un omomorfismo dalla coalgebra \(\zeta\) a se stessa; anche \(\mathrm{id}_Z\) lo è, e per unicità si conclude \(f \circ \zeta = \mathrm{id}_Z\). Per l'equazione inversa:
  \[
    \zeta \circ f = F(f) \circ F(\zeta) = F(f \circ \zeta) = F(\mathrm{id}_Z) = \mathrm{id}_{F(Z)}.
  \]
  Dunque \(f = \zeta^{-1}\) e \(\zeta\) è un isomorfismo.
\end{proof}

Il lemma dice che la coalgebra finale è un \textit{punto fisso} dell'endofuntore \(F\): si ha \(Z \cong F(Z)\). Per il funtore \(\mathrm{Seq}_A \) questo diventa \(A^\infty \cong \{\bot\} \sqcup (A \times A^\infty)\), l'equazione caratteristica delle sequenze finite e infinite \cite{jacobs2017}.

La finalità fornisce un principio di \textbf{definizione per coricorsione}: per definire una funzione \(h : X \to Z\) nella coalgebra finale è sufficiente dotare \(X\) di una struttura di \(F\)-coalgebra \(c : X \to F(X)\), e la funzione di comportamento \(\mathrm{beh}_c\) è l'unica funzione che si ottiene in questo modo. Non è necessario specificare \(h\) globalmente: è sufficiente descrivere come si comporta un passo alla volta, e la finalità garantisce l'esistenza e unicità della funzione definita.

\section{Coinduzione come principio di dimostrazione}

L'unicità del morfismo \(\mathrm{beh}_c\) non è solo una proprietà tecnica: è il fondamento del \textbf{principio di coinduzione} \cite{jacobs2017,rutten2000}. Nella sua forma più generale, esso afferma che per dimostrare l'uguaglianza di due funzioni \(f, g : X \to Z\) verso la coalgebra finale è sufficiente esibire una struttura di coalgebra \(c : X \to F(X)\) rispetto alla quale sia \(f\) che \(g\) siano omomorfismi, poiché l'unicità di \(\mathrm{beh}_c\) ne forza allora la coincidenza.

\section{Bisimulazione}

Il principio di coinduzione afferma che due stati con la stessa funzione di comportamento sono indistinguibili. La \textit{bisimulazione} è il concetto duale che cattura questa equivalenza comportamentale direttamente in termini di relazioni tra stati, senza fare riferimento alla coalgebra finale \cite{rutten2000}.

\begin{definition}[Bisimulazione]
  Siano \(c : X \to F(X)\) e \(d : Y \to F(Y)\) due \(F\)-coalgebre. Una \textbf{bisimulazione} tra \(c\) e \(d\) è una relazione \(R \subseteq X \times Y\), con proiezioni \(\pi_1 : R \to X\) e \(\pi_2 : R \to Y\), tale che esiste una struttura di \(F\)-coalgebra \(e : R \to F(R)\) che rende \(\pi_1\) e \(\pi_2\) omomorfismi di coalgebre, cioè che rende commutativo il seguente diagramma:
  \[
    \begin{tikzcd}[column sep=large, row sep=large]
      F(X) & F(R) \arrow[l, "F(\pi_1)"'] \arrow[r, "F(\pi_2)"] & F(Y) \\
      X \arrow[u, "c"] & R \arrow[u, "e"] \arrow[l, "\pi_1"] \arrow[r, "\pi_2"'] & Y \arrow[u, "d"']
    \end{tikzcd}
  \]
  ovvero tale che esista \(e : R \to F(R)\) con \(F(\pi_1) \circ e = c \circ \pi_1\) e \(F(\pi_2) \circ e = d \circ \pi_2\). Se \(X = Y\) e \(c = d\), si parla di bisimulazione sulla singola coalgebra \(c\). Questa formulazione è nota come bisimulazione \textit{alla Aczel--Mendler} \cite{jacobs2017}.
\end{definition}

La definizione afferma che \((x, y) \in R\) quando \(x\) e \(y\) si comportano allo stesso modo: ogni mossa di \(x\) può essere imitata da \(y\) restando in \(R\), e viceversa. Il diagramma commutativo esprime che la relazione \(R\), vista come coalgebra, è compatibile con entrambe le strutture di transizione tramite le proiezioni.

\begin{proposition}
  Le bisimulazioni godono delle seguenti proprietà di chiusura:
  \begin{enumerate}
    \item La relazione identità \(\Delta_X = \{(x, x) \mid x \in X\}\) è una bisimulazione su ogni coalgebra \(c : X \to F(X)\).
    \item Se \(R\) è una bisimulazione tra \(c\) e \(d\), allora la relazione inversa \(R^{-1} = \{(y, x) \mid (x, y) \in R\}\) è una bisimulazione tra \(d\) e \(c\).
    \item Se \(R\) è una bisimulazione tra \(c : X \to F(X)\) e \(d : Y \to F(Y)\), e \(S\) è una bisimulazione tra \(d\) e \(e : Z \to F(Z)\), allora la composizione relazionale \(S \circ R = \{(x, z) \mid \exists y \in Y.\, (x,y) \in R \wedge (y,z) \in S\}\) è una bisimulazione tra \(c\) e \(e\).
  \end{enumerate}
\end{proposition}

\begin{proof}
  \begin{enumerate}
    \item La struttura di coalgebra su \(\Delta_X\) è data da \(c\) stesso, poiché \(\Delta_X \cong X\) tramite entrambe le proiezioni, che coincidono con \(\mathrm{id}_X\). La commutatività dei diagrammi segue immediatamente.
    \item Se \(e : R \to F(R)\) testimonia che \(R\) è una bisimulazione, la stessa struttura \(e\) su \(R^{-1} \cong R\) testimonia che \(R^{-1}\) è una bisimulazione, scambiando i ruoli di \(\pi_1\) e \(\pi_2\).
    \item Siano \(e_R : R \to F(R)\) e \(e_S : S \to F(S)\) le strutture che testimoniano le due bisimulazioni. Su \(S \circ R\) si definisce una struttura di coalgebra usando l'azione di \(F\) sulle proiezioni: poiché \(F\) preserva le composizioni e le proiezioni di \(S \circ R\) si fattorizzano attraverso quelle di \(R\) e \(S\), la commutatività dei diagrammi richiesti segue per funtorialità.
  \end{enumerate}
\end{proof}

Le proprietà di chiusura mostrano che la \textit{bisimilarità} — la relazione \(\sim\) definita ponendo \(x \sim y\) quando esiste almeno una bisimulazione che contiene la coppia \((x, y)\) — è una relazione di equivalenza. Essa è inoltre la massima bisimulazione: l'unione di tutte le bisimulazioni è ancora una bisimulazione.

Il legame fondamentale tra bisimulazione e coinduzione è espresso dal seguente teorema.

\begin{theorem}
  Sia \(\zeta : Z \to F(Z)\) la coalgebra finale. Due stati \(x, x' \in X\) di una coalgebra \(c : X \to F(X)\) sono bisimili se e solo se hanno lo stesso comportamento:
  \[
    x \sim x' \iff \mathrm{beh}_c(x) = \mathrm{beh}_c(x').
  \]
\end{theorem}

\begin{proof}
  (\(\Rightarrow\)) Sia \(R\) una bisimulazione su \(c\) con \((x, x') \in R\). Allora \(R\) è essa stessa una coalgebra, e le proiezioni \(\pi_1, \pi_2 : R \to X\) sono omomorfismi. Quindi \(\mathrm{beh}_c \circ \pi_1\) e \(\mathrm{beh}_c \circ \pi_2\) sono entrambi omomorfismi da \(R\) a \(\zeta\). Per unicità dell'omomorfismo verso la coalgebra finale, \(\mathrm{beh}_c \circ \pi_1 = \mathrm{beh}_c \circ \pi_2\), cioè \(\mathrm{beh}_c(x) = \mathrm{beh}_c(x')\).

  (\(\Leftarrow\)) Supponiamo \(\mathrm{beh}_c(x) = \mathrm{beh}_c(x')\). La relazione \(R = \{(u, v) \in X \times X \mid \mathrm{beh}_c(u) = \mathrm{beh}_c(v)\}\) è una bisimulazione: la struttura di coalgebra su \(R\) è data da \(\langle c \circ \pi_1, c \circ \pi_2 \rangle\), e la commutatività dei diagrammi segue dal fatto che \(\mathrm{beh}_c\) è un omomorfismo. Poiché \((x, x') \in R\), si conclude \(x \sim x'\).
\end{proof}

Il teorema afferma che la bisimilarità e l'uguaglianza del comportamento osservabile sono la stessa cosa. Questo è il risultato fondamentale che permette di usare la bisimulazione come strumento pratico di dimostrazione: per mostrare che due stati si comportano allo stesso modo è sufficiente esibire una bisimulazione che li contenga, senza calcolare esplicitamente la loro funzione di comportamento \cite{jacobs2017,rutten2000}.

\section{Predicati e invarianti}

La teoria delle coalgebre fornisce anche un linguaggio per esprimere \textit{proprietà} degli stati di un sistema, non solo la loro struttura di transizione \cite{jacobs2017}. Sia \(c : X \to F(X)\) una \(F\)-coalgebra. Per un sottoinsieme \(P \subseteq X\), il \textit{predicato successore} \(\bigcirc_c P \subseteq X\) è definito come:
\[
  \bigcirc_c P = c^{-1}(F(P))
\]
dove \(F(P) \subseteq F(X)\) è l'immagine di \(P\) sotto l'azione del funtore \(F\) sull'inclusione \(P \hookrightarrow X\), e \(c^{-1}\) denota l'immagine inversa lungo \(c\). In parole: \(\bigcirc_c P\) è l'insieme degli stati la cui transizione cade interamente in \(F(P)\), ovvero degli stati i cui successori soddisfano tutti \(P\).

Per il funtore \(\mathrm{Seq}_A : X \mapsto \{\bot\} \sqcup (A \times X)\) e una coalgebra \(c : S \to \{\bot\} \sqcup (A \times S)\), si ha \(F(P) = \{\bot\} \sqcup (A \times P)\), e quindi:
\[
  \bigcirc_c P = \{x \in S \mid c(x) = \bot \text{ oppure } c(x) = (a, x') \text{ con } x' \in P\}.
\]
Usando la notazione di transizione \(x \xrightarrow{a} x'\), questa condizione si legge: \(\bigcirc_c P\) vale per \(x\) se tutti gli stati successori di \(x\), qualora esistano, soddisfano \(P\).

\begin{definition}[Invariante]
  Un predicato \(P \subseteq X\) è un \textbf{invariante} della coalgebra \(c : X \to F(X)\) se \(P \subseteq \bigcirc_c P\), ovvero se \(c(P) \subseteq F(P)\). Intuitivamente: una volta che \(P\) vale per uno stato, la sua transizione rimane in \(F(P)\), ovvero tutti i successori soddisfano ancora \(P\).
\end{definition}

\begin{proposition}
  Gli invarianti godono delle seguenti proprietà di chiusura:
  \begin{enumerate}
    \item Il predicato costante vero, cioè \(X\) stesso, è un invariante.
    \item Dati due invarianti \(P_1\) e \(P_2\), la loro intersezione \(P_1 \cap P_2\) è un invariante.
    \item Dato un invariante \(P\), anche \(\bigcirc_c P\) è un invariante.
  \end{enumerate}
\end{proposition}

\begin{proof}
  \begin{enumerate}
    \item Per ogni \(x \in X\) si ha \(c(x) \in F(X)\), quindi \(X \subseteq \bigcirc_c X\).
    \item Poiché \(F\) è un funtore preserva le inclusioni, quindi \(F(P_1 \cap P_2) \subseteq F(P_1) \cap F(P_2)\). Dunque:
          \[
            \begin{aligned}
              x \in P_1 \cap P_2 & \Longrightarrow x \in P_1 \wedge x \in P_2                                                   \\
                                 & \Longrightarrow c(x) \in F(P_1) \wedge c(x) \in F(P_2) &  & \text{(\(P_1, P_2\) invarianti)} \\
                                 & \Longrightarrow c(x) \in F(P_1) \cap F(P_2)                                                  \\
                                 & \Longrightarrow c(x) \in F(P_1 \cap P_2)                                                     \\
                                 & \Longrightarrow x \in \bigcirc_c(P_1 \cap P_2).
            \end{aligned}
          \]
    \item Bisogna verificare che \(\bigcirc_c P \subseteq \bigcirc_c \bigcirc_c P\). Sia \(x \in \bigcirc_c P\), cioè \(c(x) \in F(P)\). Poiché \(P\) è un invariante, \(c(P) \subseteq F(P)\), e quindi per funtorialità \(F(c)(F(P)) \subseteq F(F(P))\). Poiché \(P \subseteq \bigcirc_c P\) per invarianza, si ha \(F(P) \subseteq F(\bigcirc_c P)\), e quindi \(c(x) \in F(P) \subseteq F(\bigcirc_c P)\), cioè \(x \in \bigcirc_c \bigcirc_c P\).
  \end{enumerate}
\end{proof}

\section{Operatori temporali}

A partire dagli invarianti si definiscono i principali operatori temporali, in forma del tutto generale per una qualunque \(F\)-coalgebra \cite{jacobs2017}.

\begin{definition}[Operatori temporali]
  Sia \(c : X \to F(X)\) una \(F\)-coalgebra e \(P \subseteq X\) un predicato arbitrario.
  \begin{enumerate}
    \item Il predicato \(\square P \subseteq X\) (\textit{d'ora in poi} \(P\)) è il \textit{massimo invariante contenuto in} \(P\):
          \[
            \square P = \bigcup \{Q \subseteq X \mid Q \text{ è un invariante e } Q \subseteq P\}.
          \]
    \item Il predicato \(\diamond P \subseteq X\) (\textit{prima o poi} \(P\)) è definito come \(\diamond P = \neg\,\square\,\neg P\).
  \end{enumerate}
\end{definition}

Il predicato \(\square P\) vale per \(x\) quando tutti i futuri stati raggiungibili da \(x\) soddisfano \(P\). Il predicato \(\diamond P\) vale quando almeno uno stato futuro soddisfa \(P\).
% Le proprietà di \textit{sicurezza} dei sistemi sono esprimibili come \(\square P\), le proprietà di \textit{vivacità} come \(\diamond P\), e entrambe le definizioni hanno senso per qualsiasi tipo di sistema coalgebrico.

\begin{proposition}
  L'operatore \(\square\) è un operatore di interno: per ogni predicato \(P, Q \subseteq X\),
  \begin{enumerate}
    \item \(\square P \subseteq P\);
    \item \(\square P \subseteq \square\square P\);
    \item se \(P \subseteq Q\) allora \(\square P \subseteq \square Q\).
  \end{enumerate}
  Inoltre, un predicato \(P\) è un invariante se e solo se \(P = \square P\).
\end{proposition}

\begin{proof}
  \begin{enumerate}
    \item Per definizione, \(\square P\) è il massimo invariante contenuto in \(P\), quindi in particolare \(\square P \subseteq P\).
    \item Bisogna mostrare che ogni \(x \in \square P\) appartiene anche a \(\square\square P\), cioè che \(\square P\) è un invariante contenuto in \(\square P\). È contenuto in \(\square P\) banalmente. È un invariante: per la proposizione precedente, \(\square P\) è un invariante (è il massimo invariante contenuto in \(P\), e l'unione di invarianti è un invariante per chiusura sotto intersezioni arbitrarie). Poiché \(\square\square P\) è per definizione il massimo invariante contenuto in \(\square P\), si conclude \(\square P \subseteq \square\square P\).
    \item Sia \(P \subseteq Q\). Ogni invariante \(R\) con \(R \subseteq P\) soddisfa anche \(R \subseteq Q\), e quindi è tra gli invarianti contenuti in \(Q\). In particolare, \(\square P\) — che è un invariante contenuto in \(P\), e quindi in \(Q\) — è uno degli invarianti la cui unione forma \(\square Q\). Dunque \(\square P \subseteq \square Q\).
  \end{enumerate}
  Per l'ultima affermazione. Se \(P\) è un invariante, allora \(P\) è esso stesso un invariante contenuto in \(P\), quindi \(P \subseteq \square P\); d'altra parte \(\square P \subseteq P\) per il punto 1, quindi \(P = \square P\). Viceversa, se \(P = \square P\), allora \(P\) è un invariante perché \(\square P\) è per costruzione un'unione di invarianti, e la proposizione precedente garantisce che l'unione di invarianti è un invariante.
\end{proof}